% =====================================================================
% Reference contributions enumeration — structural-template invariant 8.
%
% The benchmark survey ('From Copilots to Colleagues', §1) ends its
% Introduction with this exact contributions block. The shape is:
%   - A 1-sentence framing line ("This survey makes <N> contributions
%     to the field:").
%   - An \begin{enumerate} block with one \item per contribution.
%   - Each \item starts with a 2-4-word \textbf{Bold Lead.} naming the
%     contribution category, followed by a one-sentence description in
%     concrete nouns, and ends with a parenthesised section pointer
%     (\S\,N).
%
% audit_writing.audit_structural_template's invariant-8 detector
% accepts (\S\,N), (§N), (Section N), (Sec.\ N), and (Sec.\,N).
% At least 75% of items must carry such a pointer. The benchmark
% survey hits 4/4.
%
% Replace the verbatim lead phrases below with leads derived from the
% chosen thesis's argument_steps; the structure itself must not
% deviate.
% =====================================================================

\paragraph*{Contributions}
This survey makes four contributions to the field:

\begin{enumerate}
  \item \textbf{Comprehensive Autonomy Taxonomy.} We propose and
    validate an L1--L5 taxonomy of research agent autonomy
    (analogous to SAE driving levels), providing a precise
    vocabulary for characterising system capabilities and comparing
    across architectures (\S\,2).

  \item \textbf{Systematic Architecture Analysis.} We identify and
    analyze the dominant architectural patterns---single-agent
    loops, multi-agent systems, hierarchical orchestration, and
    tool-augmented agents---with a comparative framework evaluating
    trade-offs across scalability, cost, reliability, and human
    oversight (\S\,3).

  \item \textbf{Feature-Annotated System Comparison.} We provide
    detailed analysis of 17 major systems across a six-dimensional
    feature matrix, revealing historical maturation patterns and
    identifying capability gaps (\S\,4).

  \item \textbf{Research Agenda for Open Problems.} We identify six
    fundamental challenges---cognitive loops, context limitations,
    novelty evaluation, reproducibility, safety, and cost---and
    propose concrete research directions for each, synthesizing
    insights across the surveyed systems (\S\,6).
\end{enumerate}

% =====================================================================
% Anti-patterns (audit will reject):
%   - Bare bullet list (\\itemize) instead of \\enumerate — bullets
%     don't number contributions; the reader cannot reference 'Item 3'.
%   - Items without bold leads — the skim-reader has no anchor.
%   - Items without section pointers — the contributions list reads as
%     marketing rather than as a navigation aid.
%   - Items pointing at multiple sections (e.g. "(§2, §3, §4)") — pick
%     the one section that primarily carries the contribution.
%   - "We discuss several important topics." — vague, no bold lead,
%     no §N.
% =====================================================================
